\section{Discussion}
\label{sec:discussion}

\subsection{Interpretation of Results}
\label{sec:interpretation}

\paragraph{Reconciling with the negative result of prior work.} The closest
prior study reported that RAG did not reliably help LLM-based detection and
that few-shot learning was more effective~\cite{ease2025rag}. Our findings
complement, rather than contradict, that result: our results indicate that
retrieval \emph{can} deliver large gains when the knowledge base is code-only
and retrieval is matched to the classification granularity (per-file, $k=3$,
raw-code queries). One possible explanation is that the prior study's
knowledge sources (YARA rules, security advisories) operate at a different
abstraction level than the target \texttt{setup.py}, and its single dense
mechanism could not separate retrieval design from retrieval presence. The
implication for practice is that ``RAG'' is not a single intervention: both
the corpus and the retrieval strategy need to be chosen deliberately, and our
comparison provides a controlled comparison across several points in that
design space.

\paragraph{Retrieval presence matters more than retrieval mechanism in the Llama suite.}
Across dense, sparse, and hybrid variants, malicious recall settles around
0.96--0.97 versus 0.80 with no retrieval---within this suite, the dominant
factor was whether similar code was shown at all, rather than which
similarity function retrieved it. A plausible explanation is that passages of
malicious source are frequently near-duplicates (copy-pasted droppers,
shared obfuscation snippets), which both lexical and semantic similarity can
recover.

\paragraph{Sparse retrieval is a strong, low-cost baseline.} BM25 has no
embedding model, no GPU, and no index-approximation tuning, yet yields 0.968
F1 vs.\ 0.972 for a 4B embedding model. For replication-friendly pipelines on
this benchmark, we suggest starting with BM25 and upgrading only if a failure
analysis shows semantically paraphrased attacks slipping through.

\paragraph{Contrastive retrieval offers a favorable precision/recall
trade-off.} Retrieving malicious examples \emph{only} biases the model toward
positives; adding retrieved benign look-alikes lets it weigh counterevidence.
The 0.1\% FPR of contrastive retrieval is attractive when analysts must
triage alarms, and the mechanism is orthogonal to retriever choice---a
practitioner could plausibly obtain a similar effect with BM25 over both
pools.

\paragraph{Behavior summaries reduce the signal.} The AST-based
representations tested here (behavior-based methods) discard too much
(identifiers, URL structure, file names,\dots). Their recall
regresses to no-RAG levels while their recall-adjusted F1 falls below the
baseline, which we interpret as evidence that queries should stay close to
raw code whenever the generator can accept it.

\paragraph{Jev is a bounded external comparator.} Direct Jev, with no
retrieval, exceeds the Llama no-RAG baseline but does not reach the strongest
retrieved Llama result. This does not show that retrieval is unnecessary or
that either model family dominates: the comparison changes the model, serving
provider, output interface, and prompt. The two-stage behavior-guided Jev
variant performs slightly worse than direct Jev, so the tested fixed behavior
vocabulary and 0.5 threshold do not add useful classification signal here.

\subsection{Limitations}
\label{sec:limitations}

\paragraph{Controlled single-model retrieval ablation.} The controlled
retrieval findings are for one generator (\texttt{Llama-3.1-8B-Instruct})
and one benchmark. Jev adds an externally served model-family comparator,
but we did not repeat the retrieval ablation across multiple generators.

\paragraph{Single file granularity.} Analysis is limited to
\texttt{setup.py}. Malware hidden in auxiliary files is invisible to every
method tested, which partly explains the samples no method detects.

\paragraph{Fixed retrieval settings.} We fix $k=3$ (per the reference
design); a full $k\in\{1,3,5\}$ sweep and corpus-side ablations (e.g.,
leaving retrieved documents unlabeled) remain for future work.

\paragraph{AST-based query extraction is unrefined.} Our behavior
summaries are simple import/call inventories; a richer, security-focused
static analyzer might restore some of the lost signal.

\paragraph{Jev prompt and threshold dependence.} The behavior-guided
result depends on one typed behavior vocabulary, one 0.5 threshold, and
two-stage prompting. Its typed behavior alignment is not equivalent to
source-grounded explanation faithfulness.

\subsection{Threats to Validity}
\label{sec:threats}

\paragraph{Internal.} Evidence blocks differ across methods, but all Llama
configurations share one system prompt, output schema, and decoding
parameters; the Jev comparison is not controlled in the same way.
Faithfulness scoring is an automated proxy and should be read
comparatively, and parse failures are excluded from metric denominators and
reported as coverage.

\paragraph{Construct.} Labels come from the reference dataset, and samples
with unevaluable code (9/1029) were removed, slightly privileging clean
cases. We score verdicts against labels and explanations separately against
static evidence; for Jev we report typed behavior alignment rather than
explanation faithfulness.

\paragraph{External.} The malicious population is largely typosquats from
one observation period, so findings may not transfer to other attack
families or later windows, and Jev results are tied to the recorded
provider and model version. We treat the study as a controlled comparison,
not an absolute capability claim.

\subsection{Data and Code Availability}
\label{sec:avail}

Raw per-sample outputs for the seven Llama configurations and both Jev
comparators, the evaluation scripts, request templates, and run scripts are
provided as a replication package alongside this report; the underlying dataset is the replication package of
Ibiyo et al.~\cite{ease2025rag}.
\endinput
